\section{Related Work}
\label{sec:related}

\subsection{Medical Image Translation}
Unpaired cross-modality synthesis has evolved from adversarial translation
---CycleGAN~\cite{zhu2017cyclegan} and pix2pix~\cite{isola2017pix2pix} established
the unpaired and paired regimes, followed by multi-contrast MRI specialists such
as conditional GANs~\cite{dar2019cgan} and hybrid CNN--transformers
(ResViT~\cite{dalmaz2022resvit})---to diffusion-based models. SynDiff
\cite{ozbey2023syndiff} introduced adversarial diffusion with a cyclic
consistency prior for unsupervised medical translation, and
I2I-Galip~\cite{korkmaz2025i2igalip} replaced the discriminator with a
vision--language model for efficiency. In parallel, bridges connect the two
image distributions directly rather than through an isotropic-noise bottleneck:
DDPMs~\cite{ho2020ddpm} and their score-based generalization~\cite{song2021scoresde}
provide the mathematical backbone, with score-based
NCSN~\cite{song2019ncsn}, improved schedules~\cite{nichol2021iddpm}, direct
diffusion bridges with data consistency~\cite{chung2023directbridge}, dual
implicit bridges~\cite{su2023ddib}, neural Schr{\"o}dinger
bridges~\cite{kim2024nsb}, and direct iterative inversion~\cite{delbracio2023indi}
spanning the design space. SelfRDB~\cite{arslan2025selfrdb} combines a soft-prior
bridge with self-consistent recursive sampling and reports state-of-the-art
multi-contrast translation, while explicitly leaving \emph{uncertainty
quantification} of the bridge as an open problem. Our work builds directly on
this sampler: we keep the model frozen and harvest two signals that SelfRDB
discards.

\subsection{Uncertainty Quantification in Medical Imaging}
Per-pixel uncertainty is commonly obtained by MC-Dropout~\cite{gal2016mcdropout},
deep ensembles~\cite{lakshminarayanan2017ensembles}, or test-time
augmentation~\cite{ayhan2018tta}; the epistemic/aleatoric taxonomy of
\cite{kendall2017uncertainties} frames these as sampling-based epistemic
estimates. Diffusion ensembles follow the same recipe: implicit segmentation
ensembles~\cite{wolleb2022diffusionensembles} and per-pixel standard deviations
across diffusion restorations~\cite{osuna2024microscopy} established that
inter-sample dispersion tracks difficult image regions. We inherit this insight
but differ in three ways: (i) our ensemble members are sampled from a
\emph{diffusion bridge} whose stochastic endpoints and latent variables are
native to the sampler, not injected ad hoc; (ii) we pair ensemble dispersion
with a second, complementary signal---the convergence residual---that
is unique to self-consistent recursive samplers and costs nothing; and (iii) we
fuse both into a calibrated trust map with a clinical QC gate rather than
reporting dispersion alone. Notably, concurrent bridge models actively
\emph{suppress} this stochasticity for reconstruction
\cite{song2024ccbm}; we argue it is an asset when the task is trust assessment,
even though it is a defect for point synthesis.

\subsection{Hallucination Assessment and Trustworthy Radiology AI}
Closest to our goal, the Hallucination Index~\cite{tivnan2024hallucinationindex}
quantifies fabricated content of generative reconstruction models as an
\emph{image-level} metric over repeated samples; MUMINS~\cite{oliveras2026mumins}
conditions medical image synthesis on metadata and trains an uncertainty-aware
prediction head to output a spatial uncertainty map; and the measurement-geometry
analysis of \cite{jin2026trustworthy} proves when a measurement operator can
distinguish solutions, explaining hallucination from an operator-design
perspective. These approaches either require repeated sampling at image
granularity, architectural modification with retraining, or ex-ante operator
design. TrustBridge instead derives pixel-level evidence \emph{inside} the
existing sampler, requires no architectural change or labels, and additionally
supports rule-based gating: our framework is complementary to output-guardrail
lifecycles such as RAISE~\cite{cardoso2023raise} and provides the objective,
pixel-level evidence that reporting frameworks like
AI-RADS~\cite{fink2026airads} call for, while remaining fully automatic.

\paragraph{Beyond sampling dispersion: semantic and pathology-aware signals.}
Dispersion-style evidence is blind to errors the sampler makes
\emph{consistently}---a limitation we quantify in \S\ref{sec:exp}. Three
concurrent research lines attack this from the evaluation side and mark the
natural next step for trust layers like ours. CRAFT~\cite{chung2026craft}
scores generated images against clinical criteria with a foundation-model
Clinical Alignment Score, reducing hallucination-like content; ViCTr~\cite{susladkar2025victr}
makes synthesis pathology-aware with graded severity control via EWC-regularized
LoRA adaptation; and DDWM~\cite{li2025ddwm} achieves zero-shot translation by
learning only the target-domain distribution and fusing domain-invariant
anatomy through wavelet decomposition. All three require either foundation
models, retraining, or architectural change; they are complementary to
sampler-native evidence and are discussed as future work in \S\ref{sec:discussion}.
